\section{从短 clip 到复杂动作任务}

\subsection{动作不一定靠外观猜}

讲者给出了一些非常典型的动作识别例子：举重、涂眼影、快速挥动之类的动作，往往更依赖运动而不是静态外观。视觉模型如果只盯着某一帧，很容易忽略“动作过程”。

\subsection{Temporal action localization}

分类只回答“是什么动作”，但实际视频里还常常要问“动作什么时候开始、什么时候结束”。Temporal action localization 就是在时间轴上找动作边界。这个任务把片段级分类推进到“时间上的定位 + 分类”。

\subsection{Spatio-temporal detection}

如果还想知道动作发生在画面的哪个位置，那就变成空间和时间同时定位的问题。这比普通检测更难，因为每个预测不仅有 box，还要有时间维度。

\begin{warningbox}{长视频任务的核心难点}
真正困难的不是“有没有模型”，而是“模型如何在有限算力下看完足够多的内容”。因此 clip 采样、关键片段选择和高效推理，都是视频理解里绕不开的问题。
\end{warningbox}

